\subsection*{Sprint 6: Persistencia en el entorno local}
\addcontentsline{toc}{subsection}{Sprint 6: Persistencia en el entorno local}

El sexto sprint aborda la \textbf{persistencia de datos} de \textbf{Muncher},
la pieza que separa al sistema actual de uno capaz de conservar la información
de sus usuarios. Tras los sprints anteriores, la aplicación dispone de un
front-end desplegado, de una API que lo sirve bajo el mismo dominio y de la
monitorización necesaria para observar ambos, pero la colección de recetas que
la API devuelve sigue fijada en el código. Ninguna de las funcionalidades
previstas —la creación e importación de recetas, la planificación de menús o la
generación de listas de la compra— puede construirse sin un lugar donde
guardar lo que el usuario introduce y recuperarlo en sesiones posteriores.

El trabajo necesario para incorporar la persistencia supera la capacidad de un
único sprint, por lo que se reparte entre dos. Siguiendo el criterio aplicado
al resto del proyecto de construir primero en local y desplegar después, este
sprint deja resuelta la persistencia en el entorno local y en las pruebas, y el
siguiente la lleva a la nube. La división no es únicamente una cuestión de
capacidad: los mecanismos que se establecen aquí —la forma en que la API recibe
los datos de acceso a la base de datos y la forma en que se aplican los cambios
de esquema— son los mismos que emplearán después los entornos desplegados, de
modo que llegan a la nube ya probados.

El primer objetivo es levantar una \textbf{base de datos relacional} en el
entorno local y conectar la API a ella. La composición de servicios local
incorporará su propia base de datos, lo que completa la correspondencia con la
arquitectura desplegada que describe
\hyperref[sec:entorno-local]{Entorno local}. La elección de un modelo
relacional responde a la naturaleza de los datos del dominio: recetas,
ingredientes, menús y listas de la compra son entidades fuertemente
relacionadas entre sí, cuyas restricciones de integridad conviene que garantice
el propio sistema de almacenamiento y no únicamente el código de la
aplicación. La API obtendrá los datos de acceso a la base de datos a partir de
\textbf{variables de entorno}, y no del código, por dos motivos: el mismo
código podrá conectarse sin modificaciones a la base de datos local, a la de
las pruebas o a la de cualquier entorno de la nube, y ninguna credencial
quedará registrada en el repositorio, en línea con \reqref{RNF-SE-03}.

Disponer de una base de datos no basta, sin embargo, si su estructura no puede
evolucionar de forma controlada. El esquema cambiará con cada funcionalidad
que se incorpore, y esos cambios deberán aplicarse de la misma manera en todos
los entornos, sin intervenciones manuales que puedan dejarlos en estados
distintos. Por este motivo, el sprint incluye la adopción de un \textbf{gestor
	de migraciones}: cada modificación del esquema se expresará como una migración
versionada junto al código que la requiere, de manera que la estructura de la
base de datos quede sujeta al mismo control de versiones y a la misma revisión
que el resto del proyecto, y que cualquier entorno pueda reconstruirse desde
cero aplicando la secuencia completa. El sprint no crea ninguna tabla: su
alcance se limita a dejar listos los mecanismos, y el esquema se irá definiendo
en los sprints siguientes conforme lo requieran las funcionalidades.

Por último, las pruebas de extremo a extremo dispondrán de un contenedor de
base de datos propio, distinto del que se emplea durante el desarrollo, sobre
el que se ejecutarán las migraciones al comienzo de cada ejecución, de forma
que las tablas existan y estén vacías antes de la primera prueba. Partir
siempre de un estado conocido evita que el resultado de una ejecución dependa
de los datos que dejó la anterior, y permite además comprobar en cada
ejecución que la secuencia de migraciones es capaz de construir el esquema
completo.

\subsubsection*{Objetivos iniciales}

\begin{itemize}
	\item Seleccionar el motor de base de datos relacional.
	\item Incorporar una base de datos a la composición de servicios del entorno
	      local.
	\item Conectar la API a la base de datos local, obteniendo los datos de acceso
	      a partir de variables de entorno.
	\item Adoptar un gestor de migraciones y proporcionar los comandos para
	      ejecutar las migraciones en local.
	\item Dotar a las pruebas de extremo a extremo de un contenedor de base de datos
	      propio, cuyos datos de acceso se inyecten en la API.
	\item Ejecutar las migraciones al comienzo de cada ejecución de las pruebas de
	      extremo a extremo, para que las tablas estén creadas y vacías.
\end{itemize}

\subsubsection*{Historias de usuario completadas}

\begin{center}
	\begin{tabularx}{\textwidth}{|l|X|c|}
		\hline
		\textbf{Identificador} & \textbf{Nombre}                                   & \textbf{Puntos de historia} \\
		\hline
		\usref{US-CI-09}       & Base de datos local                               & 1                           \\
		\hline
		\usref{US-CI-10}       & Conexión local de la API con la base de datos     & 1                           \\
		\hline
		\usref{US-CI-11}       & Base de datos en las pruebas de extremo a extremo & 2                           \\
		\hline
		\usref{US-BD-01}       & Gestor de migraciones                             & 2                           \\
		\hline
		\usref{US-CI-12}       & Migraciones en las pruebas de extremo a extremo   & 1                           \\
		\hline
	\end{tabularx}
\end{center}
